% docs/math/shape_sensitivities_slides.tex
% Beamer slides: the shape sensitivities of JAX-DSMC, from the Boltzmann problem and the
% windowed objective to the code that estimates them, with a term-by-term audit.
% State of main on 2026-10-02.  Build: see docs/math/README.md.
\documentclass[10pt,aspectratio=169]{beamer}
\usetheme{default}
\usecolortheme{dove}
\setbeamertemplate{navigation symbols}{}
\setbeamertemplate{footline}[frame number]
\setbeamerfont{frametitle}{size=\large}
\usepackage{amsmath,amssymb}
\usepackage{booktabs}

\newcommand{\E}{\mathbb{E}}
\newcommand{\R}{\mathbb{R}}
\newcommand{\ind}{\mathbf{1}}
\newcommand{\dd}{\mathrm{d}}
\newcommand{\code}[1]{\texttt{\small #1}}
\newcommand{\ok}{\textcolor{green!50!black}{\textbf{in}}}
\newcommand{\no}{\textcolor{red!70!black}{\textbf{missing}}}
\newcommand{\out}{\textcolor{orange!80!black}{\textbf{left out}}}
\newcommand{\Lam}{\Lambda}
\newcommand{\pth}{\mathrm{post}}

\title{Shape sensitivities in JAX-DSMC}
\subtitle{From the Boltzmann problem and a windowed objective to the gradient estimator,\\
checked term by term against the code}
\date{main, 2026-10-02}

\begin{document}

\begin{frame}
\titlepage
\end{frame}

\begin{frame}{Outline}
\begin{enumerate}
\item The problem: Boltzmann equation on a body-dependent domain, windowed drag objective
\item DSMC as a Markov chain, and every place the geometry enters it
\item The value function (adjoint) and how the code carries it on particles
\item The first-variation identity: pathwise, own-reward, score, boundary-band and
  operator terms
\item Each term as implemented (\code{jax\_dsmc/shape.py},
  \code{jax\_dsmc/ensemble\_device.py}, \code{examples/ensemble\_mesh.py})
\item Audit: what is in, what is deliberately left out, what is missing
\end{enumerate}
\vspace{1em}
Notation follows \code{docs/math/ensemble\_adjoint.tex} (``the note''); equation
labels there are cited by name (note eq.\ \code{label}).
\end{frame}

% ---------------------------------------------------------------------------------------
\section{Problem}

\begin{frame}{The kinetic problem}
Flow domain $\Omega(\theta) = \text{box} \setminus B(\theta)$, the body given by a $C^1$
level set
\[
B(\theta) = \{x : \phi(x; c(\theta)) < 0\},\qquad
\phi(x;c) = \sum_j c_j B_j(x)\quad\text{(cubic B-splines, spacing } h_\phi),
\]
with $n = \nabla\phi/|\nabla\phi|$ the normal into the flow.
Boltzmann equation for the one-particle density $f(t,x,v)$:
\[
\partial_t f + v\cdot\nabla_x f = Q(f,f)\quad\text{in } \Omega(\theta),
\]
with: inflow (drifting Maxwellian flux) at $x = 0$; outflow at $x = L$; specular
$y$-walls; and on the body, for $v\cdot n > 0$,
\[
|v\cdot n|\, f(x, v) = \int_{v'\cdot n<0} K(v' \to v;\, n)\, |v'\cdot n|\, f(x, v')\,\dd v' ,
\qquad
K = (1-a)\,\delta(v - R_n v') + a\, q_w(v\mid n),
\]
$R_n v' = v' - 2(v'\cdot n)n$, $q_w$ the diffuse half-Maxwellian at the wall temperature
(note eq.\ \code{diffuse-density}).  The squash study and \code{optc} use specular bodies ($a = 0$).

\medskip
Controls: \code{radius}, \code{centre\_x}, \code{mode2}, \code{width}, \code{squash\_y}
(chained through $c(\theta)$), or the coefficients $c$ themselves (the optimiser).
\end{frame}

\begin{frame}{The windowed objective}
Mean $x$-force on the body over the window $[t_1, t_2]$, $t_2 - t_1 = W\Delta t$:
\[
D(\theta) = \frac{1}{W\Delta t}\int_{t_1}^{t_2}\!\!\int_{\partial B}\int_{v\cdot n<0}
m\,\big(v - v^{\mathrm{out}}(v)\big)_x\, |v\cdot n|\, f\,\dd v\,\dd S\,\dd t
\]
(for a diffuse wall, $v^{\mathrm{out}}$ is averaged over $K$).  In DSMC, with $F_{\mathrm{num}}$ real molecules per simulator particle,
this becomes an expectation of a sum of \emph{per-hit rewards}:
\[
D = \E\Big[\sum_{k\in\text{window}}\ \sum_{\text{body hits } j \text{ in step } k}
r_{k,j}\Big],\qquad
r_{k,j} = \frac{F_{\mathrm{num}}}{W\Delta t}\, m\,(v^{\mathrm{in}}_{k,j} - v^{\mathrm{out}}_{k,j})_x .
\]
Code: \code{ensemble\_mesh.\_drag\_reward} (per hit, zero outside the window and off the body),
\code{objective\_value} (the same sum from the surface tallies; this is what gets reported).

\medskip
Squash study (\code{examples/mesh/optimize\_squash.py}): $J(s) = D(s) - \lambda_A\, A(s)$,
with $A(s) = \pi R^2 s$ exact; $\lambda_A$ is fixed so that the free-molecular optimum is at
$s = 0.35$.  Only $\dd D/\dd s$ needs the adjoint.

\medskip
There is no terminal or per-step state reward: \code{phi\_final} $= 0$, \code{phi\_step} unused.
\end{frame}

% ---------------------------------------------------------------------------------------
\section{DSMC}

\begin{frame}{DSMC as a Markov chain (\code{simulate.\_step})}
State $Z_k = \{(x_i, v_i, a_i)\}_{i=1}^N$ (fixed slots, alive mask).  One step, in order:
\[
Z_k \xrightarrow{\ \text{move + walls}\ \mathcal W_\theta\ } Z_k^{m}
\xrightarrow{\ \text{inject}\ } \cdot
\xrightarrow{\ \text{sort}\ } \cdot
\xrightarrow{\ \text{collide (NTC/VSS)}\ \mathcal C_{\theta}\ } Z_{k+1}.
\]
\textbf{RNG contract:} every stage draws from $\mathrm{fold\_in}(\mathrm{fold\_in}(\kappa, k), \text{stage})$, so
with the key fixed each step is a deterministic function of its input.  The record
(\code{EventRecord}) stores $Z_k^m$, the collision pairs, births, and wall events.

\medskip
\textbf{Where $\theta$ enters the forward}:
\begin{enumerate}
\item[(G1)] the move: hit times $t^h$, hit points $x^h$, normals $n$, specular
  $v^{\mathrm{out}} = R_n v^{\mathrm{in}}$, and multi-bounce chains;
\item[(G2)] the per-hit reward itself, through $v^{\mathrm{out}}(n)$;
\item[(G3)] the diffuse kernel $q_w(\cdot\mid n)$ (its normal);
\item[(G4)] which particles hit (hit set), and in which step;
\item[(G5)] NTC: the attempt count $\bar n_{\mathrm{att},c} = \tfrac12 N_c(N_c-1)\Sigma_c\,\Delta t\,F_{\mathrm{num}}/V_c$
  uses the \emph{cut-cell flow volume} $V_c(\theta)$;
\item[(G6)] inject: entrants are tracked for a random fraction of $\Delta t$ (birth-step hits).
\end{enumerate}
\end{frame}

% ---------------------------------------------------------------------------------------
\section{Adjoint}

\begin{frame}{The value function}
For the expected objective $D = \E[\sum_k R_k]$ define the value
of a particle state at the post-move instant of step $k$:
\[
\lambda_k(x, v) = \frac{\delta D}{\delta f^m_k(x,v)}
\;\approx\; \E\Big[\text{rewards collected from } t_k^m \text{ on}\ \Big|\ \text{one particle at } (x,v)\Big],
\]
one simulator particle standing for $F_{\mathrm{num}}$ molecules.
For a mean-field (Boltzmann) gas, $\lambda$ solves the adjoint equation
\[
-\partial_t\lambda - v\cdot\nabla_x\lambda = \mathcal L^*_f\lambda + (\text{reward sources}),
\]
\[
(\mathcal L^*_f\psi)(v) = \!\int\! f(w)\!\int\! q\,[\psi(v') + \psi(w') - \psi(v) - \psi(w)]\,\dd\sigma\,\dd w ,
\]
with, on a specular body, $\lambda(x, v^{\mathrm{in}}) = r(v^{\mathrm{in}}) + \lambda(x, R_n v^{\mathrm{in}})$
(the hit pays its reward and the particle continues).

\medskip
The code never solves this PDE.  It carries \textbf{one value per particle slot, at that
slot's own state}, going backward through the record (route (b) of the review; note \S3).
\end{frame}

\begin{frame}{The backward recursion as implemented}
Backward through step $k$, with $\mu$ the values at $t_{k+1}$:
\begin{enumerate}
\item \textbf{Collision pull-back} $\mu \mapsto \lambda^m_k$ (\code{\_substep}, \code{\_partner\_pass}):
$S$ substeps, $K$ virtual partners $j_\kappa$ per particle, all four values from one per-cell
fit $\Lam$:
\[
\lambda_i \mathrel{+}= \sum_\kappa w_{i\kappa}\big[\Lam(v_i') + \Lam(v_{j_\kappa}') - \Lam(v_i) - \Lam(v_{j_\kappa})\big],
\qquad
w_{i\kappa} = \frac{q(g_{i\kappa})}{\Sigma_c}\,\frac{2\,n_{\mathrm{att},c}}{N_c}\,\frac{1}{KS}
\]
(note eqs.\ \code{virtual}, \code{virtual-weight}); staged nodes move values across substep boundaries (note eq.\ \code{boundary-move}).
$(S,K) = (4,8)$ collisional, $(1,2)$ collisionless.
\item \textbf{Move pull-back} (\code{\_move\_pullback}): survivors carry their value back,
plus the rewards of the hits they took during the move:
\[
\mu^{(k)}_i = \ind[a^m_{k,i}\wedge \neg\mathrm{born}_{k,i}]\Big(\lambda^m_{k,i} + \sum_j r_{k,j,i}\Big).
\]
Newborns and particles that leave the box during the move get $0$.
\end{enumerate}
The output is $\lambda^m_k$ at every post-move state (\code{Backward.lam}, $M\times N$).  The
geometry never enters this sweep, except through the record it reads.
\end{frame}

\begin{frame}{The per-cell field interpolant (\code{ensemble\_device.field\_query})}
Nodes: alive particles of the query's cell (or the $3^D$ block of cells when pooled), in
scaled coordinates $\tilde p = (x/\Delta x,\ v/v_{\mathrm{scale}})$, $v_{\mathrm{scale}} = \sqrt{kT/m}$.
The $k$ nearest nodes ($k = 32$; $96$ pooled) get weights $\omega_n = 1/(1 + (d_n/d_k)^2)$.

\medskip
Design (``quadratic\_full'', $1 + D + 3 + 6 + D$ columns, 14 in 2-D):
\[
\phi(\delta) = \big(1,\ \delta_x,\ \delta_v,\ \{\delta_{v,a}\delta_{v,b}\}_{a\le b},\ \delta_x^{2}\big),
\qquad
\Big(\sum_n\omega_n\phi_n\phi_n^\top + \rho\bar d\,\mathrm{diag}(0,1,\dots,1)\Big)\beta
= \sum_n\omega_n\phi_n\lambda_n ,
\]
solved by an unrolled Cholesky factorisation (\code{\_spd\_solve}; both value columns at once since 2026-10-02).
\[
\Lam = \beta_0,\qquad \nabla_x\Lam = \beta_{1:1+D}/\Delta x,\qquad \nabla_v\Lam = \beta_{1+D:4+D}/v_{\mathrm{scale}} .
\]
The collision pull-back uses the value; the shape gradient uses the two slopes
(\code{node\_gradients}).
\end{frame}

% ---------------------------------------------------------------------------------------
\section{First variation}

\begin{frame}{The first-variation identity for one step}
Write one step as $Z' = G_\theta(Z, \omega)$ with randomness $\omega \sim p_\theta$, reward
$R(Z, Z')$ and value $V(Z')$ of what follows.  Then
\begin{align*}
\frac{\dd}{\dd\theta}\E\big[R + V(Z')\big]
&= \underbrace{\E\big[(\partial_{Z'}R + \nabla V)\cdot\partial_\theta G\big]}_{\text{pathwise (frozen }\omega)}
+ \underbrace{\E\big[\partial_\theta R\big|_{Z'}\big]}_{\text{own reward}}\\
&\quad+ \underbrace{\E\big[(R + V)\,\partial_\theta\log p_\theta(\omega)\big]}_{\text{score}}
+ \underbrace{\int_{\Sigma_\theta}[\![R + V]\!]\, u_\Sigma\,\dd\sigma}_{\text{boundary band}}
\end{align*}
where $\Sigma_\theta$ is the set of $\omega$ at which $G_\theta$ jumps (a hit becomes a miss, or
a hit moves from one step into the next), $[\![\cdot]\!]$ the jump across it, and $u_\Sigma$ the
speed at which $\Sigma_\theta$ moves with $\theta$.  Summed over steps, with $V = \lambda_{k}$:
\[
\frac{\dd D}{\dd\theta} = \sum_k\Big(\text{pathwise}_k + \text{own}_k + \text{score}_k + \text{band}_k\Big)
+ \underbrace{\sum_k \E\big[\partial_\theta(\text{collision operator})\text{ at fixed particles}\big]}_{\text{operator term (G5)}} .
\]
Each term below is matched to one of these pieces.  None of them differentiates the
estimator of $\lambda$: the value is a cotangent, not a function of $\theta$.
\end{frame}

\begin{frame}{The replayed wall map (\code{shape.replay}, \code{shape.wall\_map})}
Under the RNG contract the move of step $k$ is a pure function
\[
(x^\pth_k, v^\pth_k, \log p_{\varepsilon,k}, T_k, \mathcal E_k) = \mathcal W_\theta(x_k, v_k, \kappa_k)
\]
($T_k$ surface tallies, $\mathcal E_k$ traced events: $x^h, v^{\mathrm{in}}, v^{\mathrm{out}}, n, t^h$ per hit).
\begin{itemize}
\item \code{shape.replay} re-runs $\mathcal W_\theta$ from the record, step by step
  (\code{lax.scan}), forms the per-hit rewards and timing rewards, and checks the replay:
  any flipped outcome is an error; continuous error $\le 10^{-9}$ (\code{REPLAY\_TOL}).
\item \code{shape.\_gradient\_step} takes \code{jax.vjp} of $\mathcal W_\theta$ w.r.t.\ the geometry
  pytree (level-set coefficients) and pulls back \emph{all} cotangents of the step in
  one batched transpose (\code{\_pullbacks}).
\end{itemize}
Only the survivors of the step contribute: $a_k \wedge a^m_k \wedge \neg\mathrm{born}_k$.
$\theta$ derivatives inside $\mathcal W$: the hit time solves $\phi(x + t v) = 0$, so
$\partial t^h/\partial c_j = -B_j(x^h)/(v\cdot\nabla\phi(x^h))$
(bisection, then one Newton step that carries the derivative),
then $x^h = x + t^h v$, $n(x^h; c)$, $v^{\mathrm{out}} = R_n v$,
$x^\pth = x^h + v^{\mathrm{out}}(\Delta t - t^h)$, repeated for each bounce.
Grazing hits: the $1/(v\cdot\nabla\phi)$ is floored at $\epsilon|v||\nabla\phi|$ in the tangent only
(\code{levelset.regularised\_quotient}; $\epsilon$ = \code{grazing}, bias $O(\epsilon)$).
\end{frame}

% ---------------------------------------------------------------------------------------
\section{Terms}

\begin{frame}{Term 1 --- pathwise $x$, $v$ (G1)}
\[
\text{pathwise}_k = \sum_{i\in\text{surv}}
\nabla_x\lambda^m_k(x^\pth_i, v^\pth_i)\cdot\partial_\theta x^\pth_i
+ \nabla_v\lambda^m_k(x^\pth_i, v^\pth_i)\cdot\partial_\theta v^\pth_i .
\]
\begin{itemize}
\item Cotangents from \code{node\_gradients} at the record's post-move state; rows
  \code{pathwise\_x}, \code{pathwise\_v} of \code{ShapeTerms}.
\item \code{hit\_fits=True}: fits only at slots with a body event in the step.  With
  $\varepsilon = 0$, every other slot has $\partial_\theta x^\pth = \partial_\theta v^\pth = 0$
  exactly, so the result is the same to the last bit at about 1/200 of the fitting work.
\item \textbf{One-sided}: nodes restricted to $v\cdot n > 0$ for the query's last hit
  normal.  A particle just reflected is fitted from particles that also just left the wall,
  not from the ones about to hit it.
\item \textbf{Hit-ahead class} (\code{shape.hit\_ahead}): nodes also restricted to
  the query's class $\kappa = $ ``this particle hits the body at a later captured step''
  (note eq.\ \code{hit-ahead}).
\end{itemize}
\end{frame}

\begin{frame}{Term 1, $\nabla_x\lambda$: the along-flight identity}
\textbf{Collisionless}: the state $(x - s v, v)$ is the same trajectory delayed by $s$.  A delay
changes the objective only by moving hits across step boundaries, at rate $T$ = the value of
the \emph{timing rewards} $\tau$ (Term 4):
\[
\lambda(x - s v, v) = \lambda(x,v) + s\,T(x,v) + O(s^2)
\quad\Longrightarrow\quad
v\cdot\nabla_x\lambda = -T .
\]
$T$ is carried by a second backward (\code{hit\_reward} $= \tau$); since 2026-10-02 both run
as one two-column sweep (PSEUDOCODE~9.7.1).  \code{shape.along\_flight} replaces the
component of $\nabla_x\lambda$ along $v$ and keeps the fitted perpendicular part:
\[
\nabla_x\lambda \leftarrow \nabla_x\lambda - \frac{v\cdot\nabla_x\lambda + T}{|v|^2}\,v .
\]
\textbf{With collisions} the backward equation $\partial_t\lambda + v\cdot\nabla_x\lambda +
\mathcal L^*_f\lambda + r = 0$ gives $v\cdot\nabla_x\lambda = -T - \mathcal L^*_f\lambda$: the
extra flight time is spent colliding at the rate of the state.  Since 2026-10-02 the code adds
it: the virtual pull-back returns each slot's increment $\Delta t\,\mathcal L^*\lambda =
\sum_{\text{substeps}}\sum_\kappa w\,\Phi$ (\code{Backward.collision}), and
\code{gradient\_variants(collision=)} uses $T + \mathcal L^*\lambda$ per step.

\medskip
Default estimator: \code{along\_pool} (pooled $3^D$ block, $k = 96$, plus along-flight),
used by \code{optimize.py} and \code{optimize\_squash.py}.
\end{frame}

\begin{frame}{Term 2 --- the reward's own dependence: the tally term (G2)}
The drag reward $r = \frac{F_{\mathrm{num}}}{W\Delta t}\,m(v^{\mathrm{in}} - v^{\mathrm{out}})_x$ depends on $\theta$
through $v^{\mathrm{out}} = R_n v^{\mathrm{in}}$ (and, for multi-bounce chains, through $v^{\mathrm{in}}$ of later
bounces).  The surface tallies are outputs of $\mathcal W_\theta$, so the term is one more
cotangent:
\[
\text{own}_k = \frac{\partial D}{\partial T_k}\cdot\partial_\theta T_k,\qquad
\frac{\partial D}{\partial T_k} = \frac{F_{\mathrm{num}}}{W\Delta t}\,\ind[k\in\text{window}]\ \text{on the body rows' } x\text{-momentum}.
\]
Code: \code{tally\_cotangent} $\times$ \code{tally\_weight} in \code{\_gradient\_step}; row
\code{pathwise\_tally}.

\medskip
On the specular cylinder in free-molecular flow this term alone is $\dd F/\dd r = F/r$
($F\propto\int\cos^2\varphi\,\dd b = 4r/3$; the new grazing rays carry $\cos^2\varphi = 0$).
Measured: $1.27\pm0.02$ vs $F/r = 1.24$ ($\times10^{-20}$).  The $\nabla\lambda$ terms carry the
flow-field (blockage) feedback.
\end{frame}

\begin{frame}{Term 3 --- the score terms (G3)}
\textbf{Diffuse normal.}  The re-emitted velocity is \code{stop\_gradient}; its law depends
on $n$:
\[
q_w(v\mid n) = \frac{v_n}{s^2}e^{-v_n^2/2s^2}\,\frac{e^{-|v_t - u|^2/2s^2}}{2\pi s^2},\qquad
\text{score}_k = \sum_j w_{k,j}\,\nabla_n\log q_w(v^{\mathrm{out}}_{k,j}\mid n_{k,j})\cdot\partial_\theta n_{k,j}.
\]
Weight: the value from the draw onward.  With a tally cotangent (the drag runs) the hit's
own reward is pathwise through the reparameterised sample, so
$w_{k,j} = \lambda^m + \sum_{j'>j} r_{k,j'}$ (\code{to\_go.post}); otherwise
$\lambda^m + \sum_{j'\ge j} r_{k,j'}$ (\code{to\_go.pre}).  $1/v_n$ is floored at
$\epsilon|v^{\mathrm{out}}|$ (bias $\approx\epsilon/2$ of the $v_t/v_n$ part).

\medskip
\textbf{Move log-probability} ($\varepsilon$-regularised flight time, accommodation draw):
cotangent \code{to\_go.start} on $\log p_{\varepsilon,k}$.  At $\varepsilon = 0$ and a fixed
accommodation it does not depend on the geometry.

\medskip
For specular bodies (\code{optc}, squash) the whole score term is exactly zero.
\end{frame}

\begin{frame}{Term 4 --- the hit-timing band (G4, step assignment)}
A displacement of the surface shifts hit times; a hit near the end of step $k$ crosses into
step $k+1$.  If the objective changes by $\Delta J_{k,j}$ when it does, the boundary term is
(density of hit times at the step end) $\times(-\Delta J)$, estimated over a band
$\beta\Delta t$:
\[
c_{k,j} = -\frac{\Delta J_{k,j}}{\beta\Delta t}\,\ind\big[t^h_{k,j} > (1-\beta)\Delta t\big],
\qquad\text{cotangent on } t^h_{k,j}\quad(\beta = 0.2).
\]
For the per-hit drag, $\Delta J = +r$ on the window's last step and $-r$ on the step
before it (\code{\_drag\_jump}): the term lives at the two window edges.  On a steady
window these cancel in expectation (gate G3a), so \code{ensemble\_mesh} reports it as
\code{timing} and \textbf{leaves it out} of \code{gradient}.

\medskip
The same $c_{k,j}$ are the \emph{timing rewards} $\tau$ whose value $T$ the along-flight
identity uses (Term 1).
\end{frame}

\begin{frame}{Term 5 --- the silhouette band (G4, hit $\leftrightarrow$ miss)}
$\phi_{\min} = \min_s\phi(x + s\hat v)$ along the undeflected line: $<0$ hits, $>0$ misses.  A
surface displacement flips the outcome where $|\phi_{\min}| < |\delta\phi|$:
\[
\frac{\dd D}{\dd\theta}\Big|_{\mathrm{sil}} = -\E\Big[\rho(0)\,(V_{\mathrm{hit}} - V_{\mathrm{miss}})\,\frac{\partial\phi_{\min}}{\partial\theta}\Big]
\approx -\frac{1}{2\varepsilon_b}\sum_{|\phi_{\min,i}|<\varepsilon_b}(V_{\mathrm{hit},i} - V_{\mathrm{miss},i})\,B_j(x^*_i)\,\frac{\partial c_j}{\partial\theta},
\]
($\partial\phi_{\min}/\partial c_j = B_j(x^*)$ by the envelope theorem).  The realised outcome's value is the reward-to-go; the
virtual one is the fitted $\lambda$ at the virtual post-move state plus the virtual hit's
reward (\code{silhouette\_cotangent}).  Each passing is counted once per trajectory.

\medskip
Specular: $V_{\mathrm{hit}}\to V_{\mathrm{miss}}$ at grazing incidence, so the term is $O(\varepsilon_b)$; the runs
use \code{silhouette\_band = 0}.  Diffuse: $O(1)$ (plate gate $|z|\le1.6$).
\end{frame}

\begin{frame}{Assembly and the control chain}
Per seed (\code{ensemble\_mesh.shape\_terms} $\to$ \code{shape.gradient\_variants}):
\[
\widehat{\nabla_c D} = \underbrace{\text{pathwise\_x} + \text{pathwise\_v} + \text{pathwise\_tally}}_{\code{pathwise}}
+ \code{score} + \code{silhouette},
\]
with \code{timing} reported but not added.
Pipeline: forward with capture $\to$ \code{shape.replay} (rewards $r$, timing rewards $\tau$,
replay check) $\to$ two-column value backward ($\lambda$ from $r$, $T$ from $\tau$) $\to$
\code{shape.gradient} (VJP per step, all cotangents batched).

\medskip
Chain to a scalar control: \code{control\_gradient} takes \code{jax.vjp} of
$\theta\mapsto c(\theta)$ (\code{levelset\_coef}); e.g.\ $\partial c/\partial r = -1$ at every knot
for the radius.

\medskip
Optimiser (\code{optimize.py}): seed mean $g$, band mask + Gaussian smoothing ($H^1$-type),
Adam, fixed largest step $s_{\max} = 0.05$, reinitialise $\phi$ to a signed distance, restore the
area by a uniform shift (a projection onto $\oint V_n = 0$).  Squash: scalar Adam on
$J'(s) = \widehat{\dd D/\dd s} - \lambda_A\,\dd A/\dd s$.
\end{frame}

% ---------------------------------------------------------------------------------------
\section{Audit}

\begin{frame}{Audit: every $\theta$-dependence of the forward}
\footnotesize
\begin{tabular}{@{}llll@{}}
\toprule
dependence & term & code & status\\
\midrule
(G1) hit times, points, specular $v^{\mathrm{out}}$ & pathwise $x, v$ & VJP of $\mathcal W_\theta$ & \ok\\
(G2) reward through $v^{\mathrm{out}}$ & tally & \code{tally\_cotangent} & \ok\\
(G3) diffuse kernel normal & score & \code{normal\_score} & \ok{} (0 specular)\\
(G3) $\varepsilon$ flight time, accommodation & score & \code{to\_go.start} & \ok{} (0 at $\varepsilon{=}0$)\\
(G4) hit $\leftrightarrow$ miss & silhouette & \code{silhouette\_cotangent} & \ok{} (0 if specular)\\
(G4) hit crosses a step boundary & timing & \code{timing\_cotangent} & \out{} (edges)\\
(G5) cut-cell volume in NTC & operator & \code{volume\_cotangent} & \ok{} (2026-10-02)\\
(G6) entrants' birth-step hits & --- & masked out & \out{} (monitored)\\
hits by particles leaving in the step & --- & masked out & \out{} (monitored)\\
\midrule
estimator of $\nabla_x\lambda$ along $v$ & along-flight & \code{along\_flight} & \ok{} ($+\mathcal L^*\lambda$, 10-02)\\
hit-ahead class of the fits & --- & \code{hit\_ahead} & \textbf{biased}; default now off\\
\bottomrule
\end{tabular}

\medskip
The VJP itself is complete for the move: every floating output of $\mathcal W_\theta$ gets a
cotangent, and the replay is checked against the record at every step.
\end{frame}

\begin{frame}{Finding 1: the hit-ahead class is biased by construction}
\code{shape.hit\_ahead} sets $\kappa_{k,i} = \exists\,k'>k$: slot $i$ has a body hit at $k'$.  The
fits of $\nabla\lambda$ then use only nodes of the query's own class.  This is the library
default (\code{shape.gradient}, \code{ensemble\_mesh.default\_config}), so \code{optimize.py}
(the \code{optc} study) runs with it.  The squash and squashkn studies pass
\code{value\_class="none"}.

\medskip
Open since 2026-09-25: the fine $\cos2\theta$ gate reads $-2.386\pm0.005$ against CRN FD $-1.59$
($z\approx-120$).  The suspected cause was slot refills: the class follows the slot, so a departed
particle inherits the hits of the newborn that takes its slot.

\medskip
That mislabelling is real (coarse circle, seed 1, \code{scratch/shape\_audit/hit\_ahead\_refill.py}):
\begin{itemize}
\item 57\,\% of slots are refilled within the 300 captured steps;
\item 30\,\% of the hit-ahead labels come from a later occupant;
\item 16\,\% of the fitted queries (\code{hit\_fits}) are mislabelled.
\end{itemize}
Corrected class: $\kappa_k = h_{k+1}\vee(\kappa_{k+1}\wedge\neg\,\mathrm{born}_{k+1})$, a reverse
\code{lax.scan}.  The paired gate on the next slide shows that the corrected class makes the
bias worse.
\end{frame}

\begin{frame}{Finding 1: paired gate (\code{scratch/shape\_audit/gate\_class\_fix.py})}
Coarse circle, level-set body, \code{mode2} ($\cos2\theta$), 4 CRN seeds (32--35), the same
records for every row; FD $= -1.594\pm0.043$, tally term $-1.596$ ($\times10^{-20}$), so the
$\nabla\lambda$ parts should sum to FD$-$tally $= +0.01\pm0.68\times10^{-21}$.

\medskip
\footnotesize
\begin{tabular}{@{}lcccc@{}}
\toprule
fit class & $x$ part & $v$ part & gradient & gradient $-$ FD (paired)\\
 & $10^{-21}$ & $10^{-21}$ & $10^{-20}$ & $10^{-20}$\\
\midrule
no class (\code{none}) & $-0.28$ & $+0.36\pm0.15$ & $-1.587\pm0.030$ & $+0.007\pm0.071$\\
hit-ahead, current & $-0.72$ & $-6.02\pm0.34$ & $-2.269\pm0.036$ & $-0.675\pm0.057$\\
hit-ahead, stopped at the next birth & $-0.80$ & $-10.82\pm0.62$ & $-2.758\pm0.069$ & $-1.164\pm0.091$\\
\bottomrule
\end{tabular}
\normalsize

\medskip
\textbf{Refilling is not the cause.}  A purer class makes the bias \emph{larger}.  The
stratified fit estimates $\nabla\lambda_\kappa$ within each class, while
\[
\lambda = P\lambda_1 + (1-P)\lambda_0 \;\Rightarrow\;
\nabla\lambda = P\nabla\lambda_1 + (1-P)\nabla\lambda_0 + \underbrace{\nabla P\,(\lambda_1 - \lambda_0)}_{\text{dropped}},
\]
with $P(z)$ = probability of a later body hit.  A reflected particle can come back, most likely
after a specular $y$-wall bounce (not checked), so $\nabla_v P\ne0$, and $\lambda_1 - \lambda_0\approx$ the next hit's
drag.  Neither the silhouette nor the timing term carries this.  On the capsule,
$P\approx0$ after a reflection, so the class was harmless there.

\end{frame}

\begin{frame}{Finding 2: the cut-volume (operator) term (built 2026-10-02)}
\small
The NTC attempt count carries $1/V_c$, and $V_c(\theta)$ is the cut cell's flow volume.
Neither the replay nor the VJP sees the collision step (\code{with\_geometry} rebuilds normals,
not volumes).  The missing term has the same form as the kernel-parameter term (note eq.\ \code{kernel-term}),
with $\partial_\theta\log q$ replaced by $\partial_\theta\log V_c^{-1}$:
\[
\frac{\dd D}{\dd\theta}\Big|_{V} = -\sum_k\sum_{c\ \text{cut}}\frac{\partial_\theta V_c}{V_c}\;
C_{k,c},\qquad
C_{k,c} = \sum_{i\in c}\sum_\kappa \tfrac12\,w_{i\kappa}\,\Phi_{i\kappa},
\]
$C_{k,c}$ being the expected change of $D$ from the collisions of cell $c$ in step $k$, and
\[
\frac{\partial V_c}{\partial c_j} = \int_c \delta(\phi)\,B_j\,\dd x = \int_{\Gamma\cap c}\frac{B_j}{|\nabla\phi|}\,\dd s ,
\]
with the same for each sub-cell of a split cell.
\begin{itemize}
\item Collisionless (Kn $\sim 10^{18}$: radius, \code{mode2}, \code{optc}, squash): $C = 0$, so the term vanishes.
\item \textbf{squashkn (Kn $= 0.5$, M 0.32 and M 10)}: nonzero; the cluster runs omitted it.
\item \textbf{Built 2026-10-02}: $C_{k,c} = \tfrac12\sum_{i\in c}\Delta t\,(\mathcal
  L^*\lambda)_i$ accumulated in the virtual pull-back (\code{Backward.cell\_collision});
  \code{shape.volume\_weights} $g_c = -\sum_k C_{k,c}/V_c$; \code{shape.volume\_cotangent}
  differentiates a smoothed-Heaviside flow volume ($32^2$ samples per cell, half-width $h_\phi/4$,
  split cells by their region raster).  Check: $\sum_c \dd V_c/\dd s = -7.068$ against
  $-\pi R^2 = -7.069$; weighted sums match FDs of the forward's MC volumes to their noise.
\end{itemize}
\end{frame}

\begin{frame}{Finding 3: the along-flight identity drops the collision operator}
The identity used is $v\cdot\nabla_x\lambda = -T$; with collisions it is
$v\cdot\nabla_x\lambda = -T - \mathcal L^*_f\lambda$.  The note justifies the omission at
$\le 0.3\,\%$ collisions per particle-step.

\medskip
squashkn at Kn $= 0.5$ ($\lambda_{\mathrm{mfp}} = 2R\cdot0.5 = 1.5$\,m; $\bar c\approx380$\,m/s):
\begin{itemize}
\item M 0.32, $\Delta t = 10^{-4}$: $\nu\Delta t\approx \bar c\Delta t/\lambda_{\mathrm{mfp}} \approx 2.5\,\%$ per
  particle-step, about $8\times$ that bound;
\item M 10, $\Delta t = 2.5\times10^{-5}$: $\approx0.6\,\%$ in the free stream, several per cent in the shock
  layer;
\item over a cell crossing ($h = 0.5$\,m): $h/\lambda_{\mathrm{mfp}}\approx 1/3$ of a collision.
\end{itemize}
So the along-flight component of $\nabla_x\lambda$ misses an $O(h/\lambda_{\mathrm{mfp}})$ part at Kn $= 0.5$.
The pooled one-sided fit (\code{one\_pool}) has no such omission.  The four variants run in
one scan, so comparing \code{one\_pool} with \code{along\_pool} on the squashkn seeds costs
almost nothing extra.  The arbiter is the Kn $= 0.5$ FD gate (G4), which has not been run.
\end{frame}

\begin{frame}{Finding 4: hits that are not differentiated}
The replay and the backward mask non-survivors.  Two kinds of body hits therefore get
neither a reward in $\lambda$ nor a VJP:
\begin{itemize}
\item entrants' birth-step hits (inflow tracks them for $\xi\in[0,\Delta t]$);
\item hits by a particle that leaves the box in the same step.
\end{itemize}
Their share is \code{newborn\_hit\_share} $=(D_{\mathrm{tally}} - \sum r)/D_{\mathrm{tally}}$.  It is
recorded, not enforced.  It is $0$ while the body is more than $|v|\Delta t$ from the inflow and
outflow faces.  It is nonzero for the \code{optc} Adam bodies that reached those faces, which
are already flagged as not quotable.

\medskip
Recommendation: fail (or warn) when \code{newborn\_hit\_share} exceeds about $10^{-3}$.

\bigskip
\textbf{Left out on purpose:} the silhouette term for specular bodies ($O(\varepsilon_b)$);
the grazing floors (measured: $+0.8\pm0.8\,\%$ at 0.05); the ridge $\rho = 10^{-8}$.  The
timing term of the drag is \emph{in} since 2026-10-03 (Finding 6).
\end{frame}

% ---------------------------------------------------------------------------------------
\section{Bias study}

\begin{frame}{Bias study (2026-10-02/03): how the runs were made}
Small runs on the 4\,GB laptop GPU, every estimator paired with the CRN FD of the same seed
(\code{scratch/bias/}: \code{kn\_gate.py}, \code{fm\_exact.py}, \code{horizon.py},
\code{summarise.py}).
\begin{itemize}
\item Squash ellipse ($R = 1.5$, level set $h_\phi = 0.0625$, grazing 0.05), warm-up 700,
  300 captured steps, window 200, CRN FD $h = 0.05$.
\item Free-molecular: fnum$\times2$ (75 particles/cell; fnum$\times1$ no longer fits with the new
  outputs).  Kn $=0.5$: fnum$\times2$, virtual pull-back $S\times K = 2\times4$.
\item All $\nabla_x\lambda$ variants from one replay (\code{shape.gradient\_variants}), so
  their differences are paired too.
\end{itemize}
\textbf{An exact free-molecular reference that failed.}  Without collisions each path is a
deterministic function of $\theta$, so the forward-mode JVP of the replayed wall map, chained
step to step, is the exact pathwise derivative (\code{fm\_exact.py}; it reproduces the
forward's drag to $2\times10^{-16}$).  The squash box is a dispersing billiard (convex body,
specular $y$ walls): tangents grow at every body bounce, and the chained JVP is heavy-tailed
--- from step 0 it ranged from $-5.5\times$ to $+8.8\times$ the tally term over 8 seeds, and even
over the 300 captured steps from $0.89\times$ to $1.40\times$ (the true $\nabla\lambda$ parts are a
few per cent).  Useless per seed; the fits' smoothing is what makes the
ensemble adjoint usable here.
\end{frame}

\begin{frame}{Free-molecular specular squash: adjoint $-$ FD, \% of FD}
\footnotesize
\begin{tabular}{@{}lcccccc@{}}
\toprule
variant & \multicolumn{2}{c}{$s=0.35$} & \multicolumn{2}{c}{$s=0.6$} & \multicolumn{2}{c}{$s=1$}\\
 & no timing & + timing & no timing & + timing & no timing & + timing\\
\midrule
\code{one} & $-3.0\pm0.8$ & $-2.1\pm1.2$ & $+1.3\pm4.1$ & $+1.3\pm3.9$ & $-1.8\pm2.3$ & $-1.7\pm2.3$\\
\code{one\_pool} & $-19.4\pm0.7$ & $-18.4\pm1.1$ & $-5.1\pm1.3$ & $-5.1\pm1.8$ & $-2.3\pm1.5$ & $-2.3\pm1.6$\\
\code{along} & $-1.7\pm1.2$ & $-0.7\pm1.7$ & $-0.7\pm2.4$ & $-0.7\pm2.4$ & $-0.9\pm1.5$ & $-0.9\pm1.5$\\
\code{along\_pool} & $-9.9\pm1.3$ & $-9.0\pm1.7$ & $-4.5\pm1.4$ & $-4.5\pm1.9$ & $-2.1\pm1.4$ & $-2.1\pm1.6$\\
\code{one\_pool\_plane} & $-2.7\pm0.6$ & $-1.7\pm0.9$ & $-2.6\pm1.2$ & $-2.6\pm1.9$ & $-1.5\pm1.5$ & $-1.5\pm1.6$\\
\code{along\_pool\_plane} & $-1.7\pm1.1$ & $-0.7\pm1.5$ & $-2.5\pm1.4$ & $-2.5\pm1.9$ & $-1.5\pm1.4$ & $-1.4\pm1.5$\\
\bottomrule
\end{tabular}
% seeds per s: {0.35: 8, 0.6: 8, 1.0: 8}


\normalsize
\medskip
8 CRN seeds per $s$ (fnum$\times2$), $\pm$1 SE of the paired difference.  ``+ timing'' adds the
window-edge term (Finding 6; 0.2-step band at $s = 0.35, 0.6$, 10-step band at $s = 1$).
\begin{itemize}
\item \code{one\_pool}, \code{along\_pool} (the cluster default) are $-19\,\%$ / $-9\,\%$ at
  $s = 0.35$ and $-5\,\%$ at $s = 0.6$: Finding 5.
\item Restricted to the hit's side (\code{*\_plane}) each entry is within $2\sigma$; pooled
  over $s$ a residual $-1.9\pm0.7\,\%$ remains (not located: wall re-hits before the
  captured horizon, the grazing cut-off).
\item The hit-ahead class (not in the table), $s = 1$, no timing: $+4\,\%$ (\code{one},
  $z = 2.3$), $+6\,\%$ (\code{along}, $z = 6.6$), $+20\,\%$ (pooled, $z\approx10$).  Finding 1
  holds in free-molecular flow too.
\item The grazing cut-off 0.05 changes the truncated JVP by $+0.8\pm0.8\,\%$.
\end{itemize}
\end{frame}

\begin{frame}{Finding 5: the pooled fit reaches across a thin body}
\code{pool = 1} draws the $k = 96$ nearest nodes from the $3\times3$ block of cells around
the hit's cell ($1.5\times1.5$\,m).  At $s = 0.35$ the body is $2sR = 1.05$\,m thick ---
two cells --- so the block around a hit on the upper face holds the particles below the
body.  Their values have nothing to do with the query's, and the one-sided rule
($v\cdot n > 0$) does not separate them: below the body $v\cdot n_{\mathrm{top}} > 0$ means
moving \emph{towards} the lower face.

\medskip
\textbf{Fix} (\code{ensemble\_device.node\_gradients(halfspace=)}): a one-sided query also
drops the nodes more than \code{halfspace} cells behind the plane through the query normal to
its hit normal,
\[
(x_{\mathrm{node}} - x_q)\cdot n_{\mathrm{hit}} \ge -\tfrac14\,\Delta x .
\]
Variants \code{one\_pool\_plane}, \code{along\_pool\_plane} (\code{shape.VARIANTS}); the
latter is now the default of \code{optimize.py}, \code{optimize\_squash.py} and
\code{ensemble\_mesh}.  It also removes most of the $v$ part's spread ($-1.1\to-0.12\times10^{-21}$
at $s = 0.35$).
\end{frame}

\begin{frame}{Finding 6: the window-edge timing term}
The per-hit drag's window edges make $\E J$ depend on which step a hit falls in.  The band
estimator (Term 4) uses the hits within a band of each edge.  It was left out (``zero mean
on a steady window'').
\begin{itemize}
\item The window is steady: the drag's relative change over the 200 steps is $-0.1\,\%$
  (specular) and $+0.8\,\%$ (diffuse), $\pm2\,\%$ per 20-step bin.
\item 0.2-step band (one step per edge, $\sim10$ hits per seed): $+2.8\pm1.2\,\%$ at $s = 1$
  on 8 seeds, and $-35\,\%$ on one diffuse seed.  This first looked like the missing
  $-2.4\,\%$ of the old fnum$\times1$ gate; it is noise.
\item 10-step band (\code{timing\_band} $\ge 1$: whole steps per edge,
  \code{ensemble\_mesh.timing\_band\_setting}), same 8 seeds: $+0.03\pm0.25\,\%$.
\end{itemize}
\textbf{Decision}: the term is part of the windowed objective's derivative (nonzero whenever
the window is not steady), so it is now in the reported gradient, with the 10-step band as
the default (\code{optimize.py}, \code{optimize\_squash.py}, \code{ensemble\_mesh}).
\end{frame}

\begin{frame}{Diffuse squash body (acc $=1$), free-molecular, $s = 1$}
\footnotesize
\begin{tabular}{@{}lcc@{}}
\toprule
variant & \multicolumn{2}{c}{$s=1$}\\
 & no timing & + timing\\
\midrule
\code{one} & $+1.0\pm3.1$ & $+1.2\pm3.0$\\
\code{one\_pool} & $-6.0\pm3.1$ & $-5.8\pm3.0$\\
\code{along} & $-0.6\pm3.4$ & $-0.4\pm3.3$\\
\code{along\_pool} & $-3.2\pm3.4$ & $-3.1\pm3.3$\\
\code{one\_pool\_plane} & $-3.1\pm3.1$ & $-2.9\pm3.0$\\
\code{along\_pool\_plane} & $-1.6\pm3.4$ & $-1.4\pm3.3$\\
\bottomrule
\end{tabular}
% seeds per s: {1.0: 8}


\normalsize
\medskip
8 CRN seeds, fnum$\times2$, silhouette band 0.03, grazing 0.05; adjoint $-$ FD in \% of FD.
The first diffuse squash gate: the gradient is
\[
\underbrace{9.12}_{\text{tally (hybrid)}} \;+\; \underbrace{9.84}_{\text{silhouette}} \;-\;
\underbrace{0.13}_{\text{normal score}} \;+\; \nabla\lambda\text{ parts }(-0.15)
\quad[10^{-21}\,\mathrm{N/m}] ,
\]
against FD $19.0$: the silhouette term is half of it, and every variant but the cross-body
\code{one\_pool} is within its $\pm3\,\%$ noise (``+ timing'': the 10-step band, $+0.2\pm0.8\,\%$).
The 0.2-step band had a $6\,\%$ SE here --- a diffuse hit's reward does not vanish at grazing
incidence --- one seed read $-35\,\%$ (Finding 6).
\end{frame}

\begin{frame}{Finding 2 resolved: the cut-volume term, built and gated}
\textbf{Machinery gate} (\code{scratch/bias/volume\_scale.py}, Kn $=0.5$, fnum$\times4$, $S\times
K = 2\times4$, 8 CRN seeds): scale the NTC volume of a set of cells by $1+\epsilon$ --- the only
change --- and compare the FD in $\epsilon$ ($\epsilon = \pm0.2$) with the adjoint's
$\sum_{c\in\text{set}} g_c V_c$:

\medskip
\centering\small
\begin{tabular}{@{}lccc@{}}
\toprule
cells scaled & FD & adjoint & $z$\\
\midrule
all & $2.11\pm0.26$ & $2.45\pm0.08$ & $-1.2$\\
cut cells & $0.69\pm0.45$ & $0.34\pm0.01$ & $+0.8$\\
\bottomrule
\end{tabular}
\quad ($10^{-3}$\,N/m)

\raggedright
\medskip
\textbf{Geometry factor}: $\sum_c \partial V_c/\partial s$ from the smoothed VJP $=-7.068$
against $-\pi R^2 = -7.069$; random-weight sums match FDs of the forward's own MC volumes.

\medskip
\textbf{Size on the squash}: $\dd D/\dd s|_V = -4.4\times10^{-4}$ at $s = 1$ ($-2.9\,\%$ of
$\dd D/\dd s$, seed SE $0.1\,\%$).  The squashkn cluster runs lacked it.
\end{frame}

\begin{frame}{Finding 3: the along-flight identity with collisions}
Kn $=0.5$, $s = 1$, fnum$\times2$, $S\times K = 2\times4$, 8 seeds.  The fitted slope
(\code{*\_pool\_plane}: no identity) against the identity's value of its component along $v$,
paired per seed (x part, $10^{-4}$\,N/m):
\[
\text{fit} - \text{identity}(T + \mathcal L^*\lambda) = -1.3\pm0.5,\qquad
\text{fit} - \text{identity}(T) = +4.9\pm0.5 .
\]
\begin{itemize}
\item The identity without $\mathcal L^*\lambda$ (the cluster runs) is off by
  $3.2\,\%$ of $\dd D/\dd s$ ($z\approx10$); with it the fit and the identity nearly agree.
\item The residual $-0.8\,\%$ ($z = 2.5$) is within what the discrete scheme allows: NTC gives a
  particle one collision chance per step wherever it sits in its cell, so a shift along $v$ is
  a position offset, not a delay; the continuous identity holds to $O(\nu\Delta t)$.
\item All variants are within $\pm1.5\,\%$ of each other once $\mathcal L^*\lambda$ is in.
\end{itemize}
\scriptsize
\begin{tabular}{@{}lcc@{}}
\toprule
variant & \multicolumn{2}{c}{$s=1$}\\
 & no timing & + timing\\
\midrule
\code{one} & $-7.3\pm6.1$ & $-7.2\pm6.2$\\
\code{one\_pool\_plane} & $-8.0\pm6.0$ & $-7.8\pm6.1$\\
\code{along} & $-6.2\pm5.9$ & $-6.0\pm6.0$\\
\code{along\_free} & $-9.9\pm5.9$ & $-9.8\pm6.0$\\
\code{along\_pool\_plane} & $-7.2\pm6.1$ & $-7.1\pm6.2$\\
\code{along\_pool\_plane\_free} & $-11.0\pm6.1$ & $-10.8\pm6.1$\\
\bottomrule
\end{tabular}
% seeds per s: {1.0: 8}

\quad against the paired FD ($h = 0.05$, $\sim18\,\%$ per-seed noise); precise FD: next slide.
\end{frame}

\begin{frame}{Kn $=0.5$: the remaining gap and what it is not}
\textbf{Precise FD} (\code{scratch/bias/fd\_only.py}, 32 CRN seeds, fnum$\times2$): $h = 0.1$:
$0.01650\pm0.00019$, $h = 0.2$: $0.01605\pm0.00010$, Richardson $0.01665\pm0.00026$ (the
paired $h = 0.05$ FD has 18\,\% per-seed noise; $h = 0.2$ has 3.5\,\%).

\medskip
\textbf{Adjoint} (8 seeds, with volume and timing terms): \code{along} $-6.9$, \code{along\_pool\_plane}
$-7.9$, \code{one\_pool\_plane} $-8.7$, old identity $-11.6\,\%$ (each $\pm1.8$).  Tally
$+$ volume alone: $0.0164$.

\medskip
\textbf{Ruled out} (paired, same forward):
\begin{itemize}
\item horizon: 300 $\to$ 450 captured steps (fnum$\times3$, $S\times K = 2\times8$, 3 seeds):
  $+0.1\pm0.4\,\%$ for the fit variants (the identity variants $+3.5\pm3.5\,\%$);
\item virtual quadrature: $S\times K = 2\times4 \to 4\times8$: $<1\,\%$ (2 seeds);
\item fit width: $k = 48, 96, 192$ pooled: within $1\,\%$; unpooled $k = 16$ reads
  $+4$--$10\,\%$ higher (noisy);
\item one-sided backward fits next to the body (\code{backward\_normals}): $<1\,\%$;
\item timing band: $+0.15\,\%$.
\end{itemize}
\textbf{Open}: the $\nabla\lambda$ parts ($x$ $+2\,\%$, $v$ $-10\,\%$ of $\dd D/\dd s$) are the only
pieces large enough; the exact (sequential) pull-back on the same seeds decides between the
virtual quadrature and the fits / linearisation.
\end{frame}

\begin{frame}{The free-molecular $-2\,\%$: noise or bias?}
\footnotesize
Per $s$, adjoint $-$ FD (\%), 8 seeds each, fnum$\times2$; pooled with inverse-variance weights:

\medskip
\centering
\begin{tabular}{@{}lcccc@{}}
\toprule
 & $s = 0.35$ & $0.6$ & $1$ & pooled \\
\midrule
\code{along\_pool\_plane} & $-1.7$ & $-2.5$ & $-1.5$ & $-1.8\pm0.7$ ($z=-2.5$)\\
\code{one\_pool\_plane} & $-2.7$ & $-2.6$ & $-1.5$ & $-2.6\pm0.5$ ($z=-5.3$)\\
tally term alone & $-1.2$ & $-2.6$ & $-1.2$ & \\
tally alone, fnum$\times1$ (old gate) & $-2.6$ & $-2.4$ & $-2.4$ & ($\pm0.9$--$1.4$ each)\\
\bottomrule
\end{tabular}

\raggedright
\medskip
\textbf{Systematic} ($\chi^2$ across $s$ consistent with one offset; two independent data sets).
It sits in the \emph{tally} term, which is exact given the hits --- the fits are not the cause.
\begin{itemize}
\item FD curvature (central difference, $h = 0.05$, on the analytic reference): $\le0.15\,\%$, and of
  the wrong sign.
\item Grazing cut-off: removing it (same forward) moves the gradient $+0.15\pm0.25\,\%$.
\item \textbf{The horizon}: particles reflected \emph{before} the captured steps bounce off the
  specular $y$ walls ($\sim$230 steps round trip) and re-hit inside the window, and the adjoint
  cannot see their first hit's dependence on the shape.  Capturing 600 instead of 300 steps
  (same forward, fnum$\times4$, 8 seeds): $+0.8\pm0.25\,\%$ at $s = 1$, $+1.1\pm0.3\,\%$ at
  $s = 0.35$.  Most of the $-2\,\%$ is truncation, not a code error; the cluster jobs
  \code{fm\_hor\_s*} capture the whole run (no truncation at all) to close it.
\item Consequence: the box's specular walls give the free-molecular case a memory of a few
  hundred steps; a captured horizon $\ge 600$ steps (or open $y$ faces) for gates at the 1\,\% level.
\end{itemize}
\end{frame}

\begin{frame}{Audit: where could a no-collision assumption hide?}
\footnotesize
The step is move (hits) $\to$ NTC collisions.  Each term's justification, re-checked with collisions:

\medskip
\begin{tabular}{@{}p{3.2cm}p{5.6cm}p{3.6cm}@{}}
\toprule
term & assumption & with collisions\\
\midrule
pathwise $x$, $v$ & $\nabla\lambda_k$ at the post-move state; $\lambda_k$ includes step $k$'s collision & holds\\
tally, score, hybrid normal & exact reward; weights $\lambda$ + later rewards & holds\\
silhouette & straight flight to the step end & holds (collisions follow the move)\\
\textbf{timing band} & a hit changing step matters only at the window edges & \textbf{wrong}: the hit/collision order swaps --- Finding 7\\
along-flight & continuous identity $v\cdot\nabla_x\lambda = -T-\mathcal L^*\lambda$ & $O(\nu\Delta t)$ in the discrete scheme (Finding 3); fit-only variants avoid it\\
cut-cell volume & --- & built (Finding 2)\\
cell membership of a hit particle & omitted (integer) & $O(\nu\Delta t)$ below the order term\\
\code{hit\_fits} & only hit slots' post-move states depend on $\theta$ & holds\\
\bottomrule
\end{tabular}
\end{frame}

\begin{frame}{Finding 7: the order of a hit and the step's collision}
NTC collides once per step, after the move.  A body hit at the end of step $k$ that the
shape change pushes into step $k+1$ collides in step $k$ as the \emph{incoming} particle
$(x_w, v^{\mathrm{in}})$ instead of the \emph{reflected} one $(x_w, v^{\mathrm{out}})$; both futures
coincide after the next move, so the objective changes by
\[
J_c = c(x_w, v^{\mathrm{in}}) - c(x_w, v^{\mathrm{out}}),\qquad c = \Delta t\,\mathcal L^*\tilde\lambda ,
\]
at \emph{every} step boundary (the window-edge reward jump is the collisionless remainder of the
same event).  Without collisions $c \equiv 0$.  Estimator (\code{shape.order\_jump},
\code{ShapeTerms.order}): $c$ fitted from the backward's per-slot increments on the approaching
and leaving sides at the hit point, whole-step band (hit times uniform within a step),
cotangent $J_c/\Delta t$ on every body hit's time.

\medskip
Kn $=0.5$, $s = 1$, 8 seeds: $+6.3\pm4.4\times10^{-5}$, i.e.\ $+0.4\pm0.3\,\%$ of $\dd D/\dd s$; the gap
stays $-6.5$ to $-8.3\,\%$ (\code{along} to \code{one\_pool\_plane}).  Small, because the incoming/reflected collision effects
differ by a small fraction of the reward; in the gradient by default.
\end{frame}

\begin{frame}{Finding 1, final: the hit-ahead class cannot be repaired}
The class $\kappa$ = ``a body hit at a later step'' is a property of the particle's
\emph{future}.  With $\lambda = P\lambda_1 + (1-P)\lambda_0$, a fit inside the class estimates
$P\nabla\lambda_1 + (1-P)\nabla\lambda_0$ and drops $\nabla P\,(\lambda_1-\lambda_0)$.  Restoring
it needs $\nabla P$ from the same local regression of the class indicator --- and that sum is
the mixed-pool fit again, with more variance.  It is unbiased only where $P$ is constant over the
fit's support, which no case here satisfies (re-hits off the box walls; collisions; diffuse
re-emission).  Measured: $+40\,\%$ ($\cos2\theta$), $+4$--$20\,\%$ (free-molecular squash).

\medskip
\textbf{What survives of the idea}: stratify by a function of the \emph{current state}.  The
one-sided $v\cdot n$ rule and the tangent-plane half-space (Finding 5) do exactly that, and are
unbiased wherever $\lambda$ is smooth inside each piece.

\medskip
\code{value\_class="hit\_ahead"} is kept, documented as diagnostic only, and warns when used.
\end{frame}

\begin{frame}{Validation status (CRN finite differences)}
\footnotesize
\begin{tabular}{@{}lll@{}}
\toprule
case & result & note\\
\midrule
level-set circle radius & $z=-1.4$; $=$ analytic body to $10^{-4}$ per seed & free-mol.\\
sphere radius (8 seeds) & $z=-1.71$ & free-mol.\\
specular capsule width, hit-ahead fit & $1.356\pm0.011$ vs FD $1.35$, $z=+0.08$ & free-mol.\\
diffuse plate (silhouette + hybrid normal) & $1.191\pm0.068$ vs $1.17$--$1.34$, $|z|\le1.6$ & free-mol.\\
$\cos2\theta$, coarse, before the class & $z=+0.7$, $-0.2$ & free-mol.\\
$\cos2\theta$, fine, with the class & $-2.386\pm0.005$ vs $-1.59$, $z\approx-120$ & Finding 1\\
$\cos2\theta$, coarse, no class (4 seeds) & paired gradient $-$ FD $= +0.007\pm0.071$ & free-mol.\\
lane vs stencil tracker & agree to $2.3\times10^{-15}$ & \\
\midrule
squash, specular, $s = 0.35/0.6/1$ & $-0.7/-2.6/+0.4\,\%$ ($\pm2$) & free-mol.\\
squash $s = 1$, diffuse & $-1.4\pm3.3\,\%$ (10-step timing band) & free-mol.\\
cut-volume machinery (volume scaling) & $z = -1.2$ (all), $+0.8$ (cut) & Kn $0.5$\\
squash $s = 1$, Kn $=0.5$ & $-7$ to $-9\pm1.8\,\%$ & \textbf{open}\\
\bottomrule
\end{tabular}

\medskip
Free-molecular shape gradients, specular and diffuse, are FD-consistent to $\sim2\,\%$ with
the new defaults.  The collisional gradient is $7$--$9\,\%$ low after the volume and
$\mathcal L^*\lambda$ terms; the cause is not yet located (previous slide).
\end{frame}

\begin{frame}{Status after the bias study, and next steps}
\textbf{Changed on main (uncommitted)}
\begin{itemize}
\item \code{value\_class="none"} default (Finding 1); \code{along\_pool\_plane} default
  variant (Finding 5); window-edge timing term in the reported gradient (Finding 6);
\item cut-volume term in every collisional gradient (Finding 2); $\mathcal L^*\lambda$ in the
  along-flight identity (Finding 3); warning above $10^{-3}$ newborn-hit share (Finding 4);
\item \code{backward\_virtual(collision\_out=, columns\_out=)}; \code{backward\_normals} option.
\end{itemize}
\textbf{Next}
\begin{enumerate}
\item Locate the Kn $=0.5$ gap: cluster-scale gate at fnum$\times1$, $S\times K = 4\times8$
  with the FD-only arbiter at $h = 0.1, 0.2$; a velocity-local fit test (the unpooled $k = 16$
  fit read higher); the small-$v\cdot n$ structure of $\lambda$ for reflected particles.
\item Re-run the squashkn optimisations only after that gap is closed; the free-molecular
  squash and \code{optc} studies can be re-run now with the new defaults (the old runs used
  \code{along\_pool}: $-9\,\%$ at $s = 0.35$, and \code{optc} the class: $+40\,\%$).
\item Use a timing band of whole steps (\code{timing\_band} $\ge 1$) for diffuse bodies.
\end{enumerate}
\end{frame}

\end{document}
